\section{Bolha}
O algoritmo \textit{Bolha} é um dos algoritmos mais simples para o problema da ordenação. Ele consistem em percorrer várias vezes o vetor, comparando elementos adjacentes e trocando-os de posição caso estejam fora da ordem desejada. Nesse processo, os elementos flutuam para suas posições finais. Em outras palavras, ``[…]  os elementos da lista são movidos para as posições adequadas de forma contínua, assim como uma bolha move-se num líquido. Se um elemento está inicialmente numa posição i e, para que a lista fique ordenada, ele deve ocupar a posição j, então ele terá que passar por todas as posições entre i e j.'' \cite[p.8]{viana2011pesquisa}.

A implementação do algoritmo \textit{Bolha} mostrada em pseudo-código abaixo faz uso de uma flag\footnote{Em programação, uma flag é uma variável que sinaliza um evento, situação ou estado do programa. O termo "flag" é uma palavra inglesa que significa "bandeira".}, para que o processo seja interrompido no momento em que não for mais necessário fazer trocas, ou seja, caso o vetor atinja o estado de ordenado antes da última iteração, evitando, assim, iterações desnecessárias. Além disso, é utilizada a função \textit{Troca}, que serve, neste algoritmo e nos seguintes, para inverter os valores de duas posições do vetor.

\lstinputlisting[language=C]{funcoes-c/a_bolha.txt}
% \begin{algorithm}
%     \caption{Bolha}\label{alg-bubble-sort:cap}
%     \DontPrintSemicolon
%     \Entrada{um vetor $v$ de $n$ elementos}
%     \Resultado{o vetor $v$ em ordem crescente}
%     \Inicio{
%         $i \gets 0$\;
%         \Repita{$flag = 1$}{
%             $flag \gets 0$\;
%             \Para{$j = 1$ \Ate $\,n - i - 1\,$}{
%                 \Se{$v[j-1] > v[j]$}{
%                     Troca($v[j-1]$, $v[j]$$)$\;
%                     $flag \gets 1$\;
%                 }
%             }
%             $i \gets i + 1$\;
%         }
%     }
% \end{algorithm}

\subsection{Prova de correção}
O laço mais externo possui o seguinte invariante: ao final de cada iteração, $v[n-i..n-1]$ já está na ordem definitiva. Pode-se então provar que o algoritmo está correto por indução em $i$. De fato, o invariante vale após a primeira iteração, quando $i = 1$, pois o laço mais interno terá "empurrado" o maior elemento de $v$ para a última posição. Agora supondo que o invariante é válido para as primeiras $i > 0$ iterações e seguindo este mesmo raciocínio, ao final da $(i+1)$-ésima iteração, o maior elemento de $v[0..n-i-1]$ estará na posição $n-i-1$, de modo que o invariante permanece válido, pois o vetor $v[n-i-1..n-1]$ conterá os $(i+1)$ maiores elementos de $v$ em ordem. Portanto, ao final da última iteração, quando $i = n$, o vetor $v$ estará em ordem.

\subsection{Complexidade de tempo e espaço}
O ponto crítico do algoritmo \textit{Bolha} é o teste condicional \textbf{se}. Assim, para determinar a complexidade temporal deste algoritmo, basta verificar quantas vezes esse teste será executado. O melhor caso ocorre quando o vetor já está em ordem, pois o laço interno não realizará trocas e a \textit{flag} continuará igual a zero, fazendo com que o laço externo execute apenas uma iteração. Desta maneira, o teste é executado $(n-1)$ vezes, logo a complexidade é $O(n)$. Já o pior caso ocorre quando o vetor é decrescente, uma vez que o maior elemento só chegará na última posição ao final de $n$ iterações do laço externo. Com isso, a quantidade de vezes que é feito o teste é dada pela seguinte soma:
\begin{equation}\label{eq:1}
    \sum_{i=0}^{n-2} (n-i-1) = \frac{n^2 - n}{2}
\end{equation}
Implicando em uma complexidade quadrática $O(n^2)$.

A complexidade espacial é sempre $O(1)$, pois o algoritmo usa espaço adicional apenas para algumas variáveis escalares.